%% LyX 2.4.4 created this file.  For more info, see https://www.lyx.org/.
%% Do not edit unless you really know what you are doing.
\documentclass[english]{article}
\usepackage[T1]{fontenc}
\usepackage[latin9]{inputenc}
\usepackage{units}
\usepackage{enumitem}
\usepackage{amsmath}
\usepackage{amsthm}
\usepackage{graphicx}
\usepackage{geometry}
\geometry{verbose,tmargin=1in,bmargin=1in,lmargin=1in,rmargin=1in}

\makeatletter
%%%%%%%%%%%%%%%%%%%%%%%%%%%%%% Textclass specific LaTeX commands.
\numberwithin{equation}{section}
\numberwithin{figure}{section}
\newlength{\lyxlabelwidth}      % auxiliary length 

%%%%%%%%%%%%%%%%%%%%%%%%%%%%%% User specified LaTeX commands.
\usepackage{fancyhdr}
\usepackage{lastpage}
\pagestyle{fancy}
\usepackage{enumitem}
\usepackage{ifthen}
%sets math fonts to be more readable
\everymath=\expandafter{\the\everymath\displaystyle}
%\DeclareMathSizes{10}{10}{10}{10}
\usepackage{multicol}
% Added by lyx2lyx
\setlength{\parskip}{\smallskipamount}
\setlength{\parindent}{0pt}

\makeatother

\usepackage{babel}
\begin{document}
\renewcommand{\author}{Dr.\,James Ashe}

\newcommand{\classNum}{232}

\newcommand{\classSec}{C}

\newcommand{\nameType}{}

\newcommand{\examType}{Exam 4 Review}

\renewcommand{\date}{Exam 3 will be be given on December $5^{\text{th}}$, 2012} %%\currenttime, \today

%%First page's header%%%%%%%%%%%%%%%%%%%%%%
\fancypagestyle{firststyle} %%header settings for first pagr
{
	\fancyhead[L]{MTH \classNum{}(\classSec{}) \author{}\\
	\textbf{\examType{}} \date}
	\fancyhead[C]{\textbf{\nameType}}
	\setlength{\headsep}{14pt}
}
%%%%%%%%%%%%%%%%%%%%%%%%%%%%%%%%%%%%%%%%%%

%%Next page's header%%%%%%%%%%%%%%%%%%%%%%%
\setlist{nolistsep}
\lhead{\classNum{}(\classSec{})} 
\chead{\textbf{\nameType}} 
\cfoot{\thepage{}~of~\pageref{LastPage}} 
\setlength{\headsep}{5pt} 
%%%%%%%%%%%%%%%%%%%%%%%%%%%%%%%%%%%%%%%%%%%

%%Various settings; see the preamble for other settings%%%%%
%%\setenumerate[0]{leftmargin=12pt,itemindent=0pt,labelwidth=0pt}
\renewcommand{\labelenumi}{\textbf{\arabic{enumi}.}}
\renewcommand{\labelenumii}{\textbf{\alph{enumii}.}}
\thispagestyle{firststyle}
%%%%%%%%%%%%%%%%%%%%%%%%%%%%%%%%%%%%%%%%%%

\vfill{}

\subsection*{8.1 Integration by Parts.}
\begin{itemize}
\item Definite and indefinite integral of functions with$\ln x$

\begin{itemize}
\item \textbf{20, 33, 51; ex. }${\displaystyle \int\ln2x\,dx,\,\int_{1}^{e}\ln2x}\,dx$.
Note that these examples, like problem 51, require substitution.
\end{itemize}
\item You will still be expected to be able to use integration by parts

\begin{itemize}
\item \textbf{ex. ${\displaystyle \int x^{2}\ln x\,dx}$.}
\item \vfill
\end{itemize}
\end{itemize}

\subsection*{8.2 Trigonometric Integrals.}
\begin{itemize}
\item Integrals of $\sin x$ or $\cos x$ with even or odd powers. 

\begin{itemize}
\item Even powers can be solved with the half-angle formula.
\item []\textbf{9, 10; ex. ${\displaystyle \int\cos^{4}2x}\,dx$}.\textbf{ }
\item Odd powers can be solved by 1) splitting off a single power, e.g.
$\sin^{5}x=\sin^{4}x\sin x$, 2) the Pythagorean identity, 3) substitution.

\begin{itemize}
\item []\textbf{11, 12; ex.${\displaystyle \int\sin^{5}x}\,dx$}. 
\end{itemize}
\end{itemize}
\item Integrals of products of $\sin x$ and $\cos x$ with even and/or
odd power.

\begin{itemize}
\item Problems with \emph{only }even powers can be solved with the half-angle
formula. 

\begin{itemize}
\item []\textbf{13, 17; ex. ${\displaystyle \int\sin^{4}x\cos^{2}x}\,dx$}.
\end{itemize}
\item Problems with at \emph{least one }odd power can be solved by 1) splitting
off a single power, e.g. $\sin^{5}x=\sin^{4}x\sin x$, 2) the Pythagorean
identity, 3) substitution. 

\begin{itemize}
\item []\textbf{15, 14; ex. ${\displaystyle \int\cos^{3}x\sin^{2}x}\,dx$.}\vfill
\end{itemize}
\end{itemize}
\end{itemize}

\subsection*{8.3 Trigonometric Substitutions}

Integrals involving $a^{2}-x^{2}$. The important substitutions are:
$x=a\sin\theta$ and $\sqrt{a^{2}-x^{2}}=a\cos\theta$. You may also
need the following reference triangle:\\
\includegraphics[width=3cm]{D:/home/jim/JamesAshe/JCSU/Calc_2_MTH_232/images/Trig_Sub_reference_Triangle_a^2-x^2.png}
\begin{itemize}
\item Definite integrals involving $a^{2}-x^{2}$. 

\begin{itemize}
\item \textbf{7-10; ex. }$\int_{0}^{3}\sqrt{36-x^{2}}\,dx$.
\end{itemize}
\item REMEMBER that the limits of integration need to changed in these problems
: $x=a\sin\theta$$\Rightarrow$$\sin^{-1}\left(\frac{x}{a}\right)=\mbox{new endpoint}$
(above $a=3$, so $\frac{3}{2}=3\sin\theta$$\Rightarrow$$\sin^{-1}\left(\frac{\nicefrac{3}{2}}{3}\right)=\frac{\pi}{6}$).
\item Indefinite integrals involving $a^{2}-x^{2}$.

\begin{itemize}
\item \textbf{11-17; ex.${\displaystyle \int\frac{dx}{(9-x^{2})^{\nicefrac{1}{2}}}}$}.
\end{itemize}
\item REMEMBER that the final result must be in terms of $x$ (there should
be no $\theta$'s left).
\end{itemize}

\end{document}
